\section{Problem Definition}

Antimicrobial resistance prediction from genomic data is the task of estimating a laboratory resistance phenotype from DNA-derived evidence. The scientific goal of this study is not to build a workflow for one specific bacterial species, but to develop a determinant-based genotype-to-phenotype prediction approach that can in principle be applied to any bacterial species for which sufficiently rich curated genotype and phenotype data are available. In this study, \textit{Salmonella enterica} is used as the case study because it is heavily represented in surveillance datasets and provides enough paired genotype and phenotype records for a credible antibiotic-specific evaluation.

At the isolate level, the input consists of curated AMR determinants derived from genomic analysis, including acquired genes and chromosomal mutations. The output is a binary antibiotic-specific phenotype, namely \textit{resistant} or \textit{susceptible}, derived from antimicrobial susceptibility testing (AST). The practical question is therefore straightforward: given the determinant profile of one isolate, can the final AST phenotype for a target antibiotic be predicted reliably?

The difficulty of the problem lies in the gap between genotype and phenotype. Some resistance mechanisms are directly associated with one antibiotic, whereas others act more broadly across a drug class. In addition, a detected determinant may co-occur with other resistance signals without being the direct cause of the target phenotype. As a result, the main challenge is not only to train a classifier, but also to define a biologically meaningful representation of genotype evidence before prediction.

This leads to the central research question of the study:

\begin{quote}
Can curated AMR determinants reliably predict antibiotic resistance phenotype, and how does prediction quality change when genotype evidence is connected to phenotype under different levels of biological constraint?
\end{quote}

The project is therefore not just a comparison of predictive models. It is a study of how curated genomic evidence can be transformed into an antibiotic-specific AMR prediction workflow that remains biologically interpretable while still achieving strong predictive performance. The \textit{Salmonella} cohort provides the concrete experimental setting in this report because the data are available at sufficient scale, but the intended scope of the method is broader than one organism. In principle, the same workflow can be transferred to other bacterial species whenever comparable curated determinant annotations and AST labels are available.
